\documentclass[norsk,a4paper,12pt]{article}
%\renewcommand{\thesection}{\Roman{section}} 
\usepackage[utf8]{inputenc}
\usepackage{graphicx} %for å inkludere grafikk
\usepackage{verbatim} %for å inkludere filer med tegn LaTeX ikke liker
\usepackage{tabularx}
\usepackage{booktabs}
\usepackage{amsmath}
\usepackage{float}
\usepackage{color}
\usepackage{listings}
\usepackage{hyperref}

\lstset{language=c++}
\lstset{basicstyle=\small}
\lstset{backgroundcolor=\color{white}}
\lstset{frame=single}
\lstset{stringstyle=\ttfamily}
\lstset{keywordstyle=\color{red}\bfseries}
\lstset{commentstyle=\itshape\color{blue}}
\lstset{showspaces=false}
\lstset{showstringspaces=false}
\lstset{showtabs=false}
\lstset{breaklines}
\lstset{postbreak=\raisebox{0ex}[0ex][0ex]{\ensuremath{\color{red}\hookrightarrow\space}}}
\usepackage{titlesec}

\setcounter{secnumdepth}{4}

\titleformat{\paragraph}
{\normalfont\normalsize\bfseries}{\theparagraph}{1em}{}
\titlespacing*{\paragraph}
{0pt}{3.25ex plus 1ex minus .2ex}{1.5ex plus .2ex}

\begin{document}

\begin{titlepage}
	\centering
	{\scshape\huge FYS4130 - Statistical Mechanics \par}
	\vspace{1cm}
	{\scshape\Large University of Oslo\par}
	\vspace{1.5cm}
	\includegraphics[width=0.25\textwidth]{uio.png}\par\vspace{1cm}
	{\Large\bfseries Obligatory assignment\par}
	\vspace{1.5cm}
	{\large\itshape Even Marius Nordhagen\par}
	\vfill
	{\large \today\par}
\end{titlepage}


\section{CLUSTER SIZE}
This problem aims an one dimensional lattice with $N$ sites at a given temperature. There is a atom at each site that can either have energy $+\epsilon$ or $-\epsilon$, and a chain of sites with energy $+\epsilon$ is called a cluster.

\subsection{Onebody probabilities}
Since each atom can be in two possible energy states, the one-particle partition function has two terms
\begin{equation}
Z_1=\sum_ie^{-\beta\epsilon_i}=e^{-\beta\epsilon}+e^{\beta\epsilon}
\end{equation}
and the probability of an atom having energy $+\epsilon$ and $-\epsilon$ is
\begin{equation}
P_+=\frac{e^{-\beta\epsilon}}{Z_1}=\frac{1}{1+e^{2\beta\epsilon}}\quad\text{and}\quad P_-=\frac{e^{\beta\epsilon}}{Z_1}=\frac{1}{1+e^{-2\beta\epsilon}}
\end{equation}
respectively.

\subsection{Probability of finding site in cluster}

\newpage
\section{ELECTRON GAS IN A MAGNETIC FIELD}
Here we study a ...

\subsection{}

\newpage
\section{QUANTUM COULOMB GAS}
We will look at a relativistic quantum gas of positive and negative charged particles with coulomb interaction. In addition we have an external box potential with walls of length $L$. Although we cannot solve this analytically, we will get some results. The given Hamiltonian is as follows
\begin{equation}
{\cal H}=\sum_{i=1}^{2N}c|\boldsymbol{p}_i|+\sum_{i<j}^{2N}\frac{e_ie_j}{|\boldsymbol{r}_i - \boldsymbol{r}_j|}
\end{equation}

\subsection{Schrödinger equation}
The general Schrödinger equation for a many-body system is 
\begin{equation}
H\Psi_n=E_n\Psi_n
\end{equation}
where $\Psi_n$ is the total wave function of state $n$ and $E_n$ is the corresponding total energy. The wave functions are position dependent and the energies are dependent on the box size,
\begin{equation}
{\cal H}\Psi_n(\{\boldsymbol{r}_i\})=\epsilon_n(L)\Psi_n(\{\boldsymbol{r}_i\}).
\end{equation}
The exact wave functions are unknown, but there are some constraints that apply in general. Since the particles are indistinguishable, all observable should be the same although we swap two coordinates. This results in
\begin{equation}
P(a,b)=P(b,a)\quad\Rightarrow\quad |\Psi(a,b)|^2=|\Psi(b,a)|^2
\end{equation}
and
\begin{equation}
\Psi(a,b)=e^{i\phi}\Psi(b,a).
\end{equation}
Swapping twice must give back the initial wave function, so there are two possible choices of $\phi$: 0 and $2\pi$. The first one gives a symmetric total wavefunction under exchange of two particles, which is antisymmetric for the second choice. We denote the former as bosons and the latter as fermions, where the Pauli principle is a consequence of antisymmetry.

\subsection{Scaling}
We will now scale the coordinates with respect to $L$ such that the box gets size $1^d$ where $d$ is the number of dimensions, matematically written as
\begin{equation}
\boldsymbol{r}_i'=\boldsymbol{r}_i/L.
\end{equation}
Using that the operator of linear momentum is given by 
\begin{equation}
\boldsymbol{p}_i=-i\hbar\nabla_i=-i\hbar\Big(\frac{\partial}{\partial x_1},\frac{\partial}{\partial x_2},\hdots,\frac{\partial}{\partial x_d}\Big),
\end{equation}
the corresponding scaled momentum operator is 
\begin{equation}
\boldsymbol{p}_i'=-i\hbar\Big(\frac{\partial}{\partial x_1'},\frac{\partial}{\partial x_2'},\hdots,\frac{\partial}{\partial x_d'}\Big)= \frac{\boldsymbol{p}_i}{L}
\end{equation}
which is the scaling of the first term of the Hamiltonian. The second term is scaled similarly, and by denoting the scaled Hamiltonian as ${\cal H}'$, we get ${\cal H}={\cal H}'/L$. We are now set to scale the Schrödinger equation, which becomes
\begin{equation}
\frac{{\cal H}'}{L}\Psi_n(\{\boldsymbol{r}_i\})=\frac{\epsilon_n(1)}{L}\Psi_n(\{\boldsymbol{r}_i\})
\end{equation}
where the scaled energy is a function of 1 since the length scale now is 1, such that $\epsilon_n(L)=\epsilon_n(1)/L$.

\subsection{Collapse V + T}
We have a fixed number of particles and the energy cna fluctuate, so we are in the canonical ensemble. The one particle partition function is given by
\begin{equation}
Z_1=\sum_{n=0}^{\infty}e^{-\beta\epsilon_n(L)}
\end{equation}
which in three dimensions is given by
\begin{align}
Z_1(V,T)&=\sum_{n_x,n_y,n_z=0}^{\infty}\exp{\big(-\beta\epsilon_{n_x}(L)\epsilon_{n_y}(L)\epsilon_{n_z}(L)}\big)\\
&=\sum_{n_x,n_y,n_z=0}^{\infty}\exp\Big({-\frac{\epsilon_{n_x}(1)\epsilon_{n_y}(1)\epsilon_{n_z}(1)}{k^3T^3L^3}}\Big).
\end{align}
Since $L^3=V$, the one particle partition function is a function of $VT^3$. The total partition function is given by 
\begin{equation}
Z(N,V,T)=\frac{Z_1(V,T)^N}{N!}=\frac{1}{N!}\sum_{n_{\boldsymbol{k}}=0}^{\infty}e^{-N\frac{\epsilon_{n_{\boldsymbol{k}}}(1)}{k^3VT^3}}={\cal Z}(N,VT^3)
\end{equation}
since the particles are indistinguishable.

\subsection{Energy and pressure relation}
The internal energy in this system is a sum over all energy levels times the probability of that level. Further we observe that we can differentiate the exponential with respect to $\beta$ to obtain the same expression,
\begin{equation}
E=\frac{1}{Z}\sum_n\epsilon_ne^{-\beta\epsilon_n}=-\frac{1}{Z}\sum_n\frac{\partial}{\partial\beta}e^{-\beta\epsilon_n}=-\frac{1}{Z}\frac{\partial Z}{\partial\beta}.
\end{equation}
The pressure is given by the differential
\begin{equation}
P(N,V)=-\Big(\frac{\partial F}{\partial V}\Big)_N
\end{equation}
where $F$ is Helmholtz free energy,
\begin{equation}
F=-kT\ln Z.
\end{equation}
We therefore end up with the expressions 
\begin{align}
E&=-\frac{1}{Z}\frac{\partial Z}{\partial\beta}\\
P&=kT\frac{1}{Z}\frac{\partial Z}{\partial V}
\end{align}
such that they are related through the partition function. We use the same partition function as we found in the previous exercise, and obtain
\begin{align}
\frac{\partial Z}{\partial \beta}&=-3\frac{\beta^2\epsilon_{n_{\boldsymbol{k}}}(1)}{V}Z\\
\frac{\partial Z}{\partial V}&=\frac{\beta^3\epsilon_{n_{\boldsymbol{k}}}(1)}{V^2}Z.
\end{align}
The energy then reads
\begin{equation}
E=3\frac{\beta^2\epsilon_{n_{\boldsymbol{k}}}(1)}{V}=3kT\frac{\beta^3\epsilon_{n_{\boldsymbol{k}}}(1)}{V^2}V=3PV.
\end{equation}

Notice that the calculations are done for the one particle partition function, but it is easy to imagine that the full partition function will give the same result. 

\subsection{Non-relativistic gas}
We now switch from a relativistic gas to a non-relativistic one and redo the exercises in four dimensions. The Hamiltonian is then given by
\begin{equation}
{\cal H}=\sum_{i=1}^{2N}\frac{\boldsymbol{p}_i^2}{2m}+\sum_{i<j}^{2N}\frac{e_ie_j}{|\boldsymbol{r}_i - \boldsymbol{r}_j|}.
\end{equation}
We then do a change of scale similarly as in exercise 3.2, and get the Schrödinger equation 
\begin{equation}
\frac{{\cal H}'}{L^2}\Psi_n(\{\boldsymbol{r}_i\})=\frac{\epsilon_n(1)}{L^2}\Psi_n(\{\boldsymbol{r}_i\})
\end{equation}
where  ${\cal H}'$ is the scaled Hamiltonian. Further the partition function is 
\begin{align}
Z_1&=\sum_{n_x,n_y,n_z,n_u=0}^{\infty}\exp\Big({-\beta^4\epsilon_{n_x}(L)\epsilon_{n_y}(L)\epsilon_{n_z}(L)\epsilon_{n_u}(L)}\Big)\\
&=\sum_{n_x,n_y,n_z,n_u=0}^{\infty}\exp\Big({-\frac{\beta^4\epsilon_{n_x}(1)\epsilon_{n_y}(1)\epsilon_{n_z}(1)\epsilon_{n_u}(1)}{L^8}}\Big)
\end{align}
and since we are in 4D, the 4-volume is $V_4=L^4$. We can reuse the expressions for energy and pressure from above, and find a relation between the energy and pressure as we did in exercise 3.4. Again we calculate the differentials
\begin{align}
\frac{\partial Z}{\partial \beta}&=-4\frac{\beta^3\epsilon_{n_{\boldsymbol{k}}}(1)}{V_4^2}Z\\
\frac{\partial Z}{\partial V_4}&=2\frac{\beta^4\epsilon_{n_{\boldsymbol{k}}}(1)}{V_4^3}Z.
\end{align}
and find the relation
\begin{equation}
E=4\frac{\beta^3\epsilon_{n_{\boldsymbol{k}}}(1)}{V_4^2}=4kT\frac{\beta^4\epsilon_{n_{\boldsymbol{k}}}(1)}{V_4^3}=2PV
\end{equation}.

Again this is results is found from the one particle partition function, but it is general for a non-relativistic interacting quantum gas in 4D.

\newpage
\section{NON-INTERACTING FERMIONS}
In this problem we study a non-interacting fermi gas with chemical potential $\mu$. The particle states are given by the vector $\vec{k}$, which is $d$-dimensional where $d$ is number of spatial dimensions, and the corresponding energy is $\epsilon(\vec{k})$.

\subsection{Joint probability}
The joint probability of finding a particle in a state given by the occupation numbers $\{n_{\vec{k}}\}$ is given by the general probability statement in a grand canonical ensemble
\begin{equation}
P(\{n_{k}\})=\frac{e^{-\beta n_k(\epsilon_k-\mu)}}{\Xi}
\end{equation}
where $\Xi$ is the grand canonical partition function, given by
\begin{equation}
\Xi
\end{equation}

\subsection{}

\subsection{Maximum entropy of random variable}
A random variable has $l$ discrete outcomes with their own respective probabilities. Since the probabilities are independent, we can find the information entropy
\begin{equation}
S=H(p_1,\hdots,p_n)=-\sum_{n=1}^lp_n\log p_n
\end{equation}
where $p_i$ is the probability of outcome $i$. We cannot find the exact entropy without knowing the probabilities, but we can find which probabilities that give maximum entropy.

\subsection{Entropy of fermi gas}

\iffalse
\begin{figure}[!htbp]
\centering
\includegraphics[width=100mm]{example}
\caption{Time evolution of the motion in y-direction, where both the numerical and classical solution are plotted. Plotted over 6 periods and 600 timesteps per period. \label{y_time}}
\end{figure}
\fi
\end{document}
